
\documentclass{article}
%%%%%%%%%%%%%%%%%%%%%%%%%%%%%%%%%%%%%%%%%%%%%%%%%%%%%%%%%%%%%%%%%%%%%%%%%%%%%%%%%%%%%%%%%%%%%%%%%%%%%%%%%%%%%%%%%%%%%%%%%%%%%%%%%%%%%%%%%%%%%%%%%%%%%%%%%%%%%%%%%%%%%%%%%%%%%%%%%%%%%%%%%%%%%%%%%%%%%%%%%%%%%%%%%%%%%%%%%%%%%%%%%%%%%%%%%%%%%%%%%%%%%%%%%%%%
%TCIDATA{OutputFilter=LATEX.DLL}
%TCIDATA{Version=5.50.0.2953}
%TCIDATA{<META NAME="SaveForMode" CONTENT="1">}
%TCIDATA{BibliographyScheme=Manual}
%TCIDATA{Created=Friday, August 07, 2026 13:31:29}
%TCIDATA{LastRevised=Tuesday, October 27, 2026 16:32:16}
%TCIDATA{<META NAME="GraphicsSave" CONTENT="32">}
%TCIDATA{<META NAME="DocumentShell" CONTENT="Standard LaTeX\Blank - Standard LaTeX Article">}
%TCIDATA{CSTFile=40 LaTeX article.cst}

\newtheorem{theorem}{Theorem}
\newtheorem{acknowledgement}[theorem]{Acknowledgement}
\newtheorem{algorithm}[theorem]{Algorithm}
\newtheorem{axiom}[theorem]{Axiom}
\newtheorem{case}[theorem]{Case}
\newtheorem{claim}[theorem]{Claim}
\newtheorem{conclusion}[theorem]{Conclusion}
\newtheorem{condition}[theorem]{Condition}
\newtheorem{conjecture}[theorem]{Conjecture}
\newtheorem{corollary}[theorem]{Corollary}
\newtheorem{criterion}[theorem]{Criterion}
\newtheorem{definition}[theorem]{Definition}
\newtheorem{example}[theorem]{Example}
\newtheorem{exercise}[theorem]{Exercise}
\newtheorem{lemma}[theorem]{Lemma}
\newtheorem{notation}[theorem]{Notation}
\newtheorem{problem}[theorem]{Problem}
\newtheorem{proposition}[theorem]{Proposition}
\newtheorem{remark}[theorem]{Remark}
\newtheorem{solution}[theorem]{Solution}
\newtheorem{summary}[theorem]{Summary}
\newenvironment{proof}[1][Proof]{\noindent\textbf{#1.} }{\ \rule{0.5em}{0.5em}}
\input{tcilatex}
\begin{document}


\begin{enumerate}
\item[3] a) Considere una carga puntual $q=-5$ [$\mu C$] que se ubica en una
regi\'{o}n del espacio donde hay un campo magn\'{e}tico de intensidad $B=0,02
$ [T]. Si la rapidez de la carga puntual y la direcci\'{o}n del campo magn%
\'{e}tico se indican en la Fig. 3a, determine la fuerza magn\'{e}tica $\vec{F%
}_{m}$ ejercida sobre la carga puntual.%
\[
fig3a
\]

\item[3] b) Considere un lazo de corriente en forma rectangular de lados $p=5
$ [cm] y $q=8$ [cm] que se encuentra sobre el plano $xy$, como muestra la
Fig. 3b. Por el lazo de corriente circula un corriente en sentido
antihorario $I=2,5$ [A] y en la regi\'{o}n hay un campo magn\'{e}tico $\vec{B%
}=0,8j-0,3\hat{k}$ [T]. Determine:

\begin{enumerate}
\item La fuerza magn\'{e}tica  $\vec{F}_{m}$ sobre el tramo $12$ y $23$ del
lazo de corriente.

\item El momento magn\'{e}tico $\vec{\mu}$ del lazo de corriente.

\item El torque magn\'{e}tico $\vec{\tau}$ que ejerce el campo sobre la
espira.
\end{enumerate}
\end{enumerate}

\bigskip

\end{document}
